\section{Experimental Setup}
\label{sec:setup}

\subsection{Dataset requirements}
\label{sec:setup-requirements}
Before choosing datasets, we agreed on the following requirements:
\begin{description}
  \item[R1] A real social network with human-like structure: groups form through
            reachability and shared interests. The topic does not matter, because the
            message we propagate is synthetic.
  \item[R2] A node is any identity that can pass a message on.
  \item[R3] Every node has metadata that reflects its thoughts or interests (text,
            interests, topics). This metadata is used only as features.
  \item[R4] Nodes without metadata are kept (zero features plus a flag), because
            dropping them would delete edges.
  \item[R5] Any social tie counts as an edge, and direction is not required.
  \item[R6, R7] Neither labels nor timestamps are required. Success means spread is
            contained in simulation, and spread probabilities come from a model.
  \item[R8] The data is available offline, with no API hydration needed.
  \item[R9] Start small (up to about 100k nodes) and scale up later.
  \item[R10] At least two datasets.
\end{description}
We checked eight candidate datasets against these requirements using the data files
themselves, not their documentation: TIMME, Pokec, Gab, Reddit hyperlinks,
VoterFraud2020, Orkut, SNAP Facebook-ego and SNAP Twitter-ego. TIMME and Pokec are
used. VoterFraud2020 failed: its retweet table points
to tweets, not users, so no user-to-user graph can be built, and it has no per-user
text. Orkut failed because its only metadata is group membership, which is also its
ground truth. SNAP Facebook-ego failed on size and on its metadata being demographic.
SNAP Twitter-ego passed but is a union of ego networks, so its bridges would mostly be
the egos themselves; it is kept only as a backup. Gab and Reddit hyperlinks are held in
reserve.

\subsection{Datasets}
\label{sec:setup-data}
\begin{table}[t]
\centering
\caption{Datasets after preprocessing (largest connected component). Edge counts are
undirected unless marked $^{\dagger}$. Counts come from the project's loader logs.}
\label{tab:datasets}
\small
\begin{tabularx}{\linewidth}{@{}l r r r >{\raggedright\arraybackslash}X@{}}
\toprule
Dataset & Nodes & Edges & Mean deg. & Node text \\
\midrule
TIMME \texttt{P\_all}          & 20,811    & 4,106,446  & 394.6 & bio + latest tweet (20,401 with text); GloVe tweet average \\
Pokec slice (Kysuck\'e Nov\'e Mesto) & 5,599 & 45,043 & 16.1 & free-text interest fields (4,252 with text) \\
Pokec full                    & 1,632,803 & 30,622,564$^{\dagger}$ & --- & free-text interest fields \\
\bottomrule
\end{tabularx}

\smallskip
\raggedright\footnotesize
$^{\dagger}$Directed friendship rows as distributed (54.3\% reciprocated). The full
Pokec graph is currently used only to characterise community structure. All Pokec
experiments in \cref{sec:setup} use the district slice.
\end{table}


\paragraph{TIMME (\texttt{P\_all})~\citep{xiao2020timme}.}
A 2019 crawl of US political Twitter, centred on politicians: it contains all
politicians, their followers and followees with at least 20 links, and 3,000 randomly
chosen users who discuss political issues. It has five directed relation types (follow
915,438; like 1,794,111; mention 1,571,937; retweet 1,684,023; reply 530,598). These
are raw directed rows, 6,496,107 in total. Pooling them, making every edge two-way,
removing duplicates and keeping the largest connected component gives one undirected
graph with 4,106,446 edges (\cref{tab:datasets}). Text: the user's bio (84.4\% non-empty) and
latest tweet (93.4\%) from the shipped user-info file, plus the shipped 300-d GloVe
average over each user's tweets~\citep{pennington2014glove}. Party labels come from
two files. The politician file lists 585 politicians (313 R, 270 D, 2 independents);
581 of the 583 R/D politicians are in \texttt{P\_all}. The additional-label file has
633 users in \texttt{P\_all}, of whom 628 are R/D and 5 independent. The two files do
not overlap, so $581+628=1{,}209$ users have an R/D label. Three of them fall outside
the largest connected component, which leaves 1,206 labelled users (644 D, 562 R).
Independents are treated as unlabelled. The labels are used
\emph{only} for the party-misfit check (\cref{sec:method-detect}) and never for model
selection. Caveat: the crawl is centred on seed accounts, so some apparent bridges may
be artefacts of how the data was collected (\cref{sec:method-limits}).

\paragraph{Pokec~\citep{takac2012pokec,leskovec2014snap}.}
A Slovak social network from 2012, with 1,632,803 users and 30,622,564 directed
friendship links, distributed through SNAP. It has no timestamps. Profiles contain
Slovak free text: hobbies, music, movies, books, sports and so on. 60.7\% of users fill
in at least one interest field. Development and all current experiments use a
district-level slice (Kysuck\'e Nov\'e Mesto, \v{Z}ilina region). The content view uses
only the interest text. Region is not used as content, because it is constant inside
the slice and would otherwise mostly encode geography. Pokec has no labels, so it is
evaluated only through spread simulation.

\subsection{Text features}
\label{sec:setup-features}
All text is encoded with one multilingual sentence encoder,
\texttt{paraphrase-\allowbreak multilingual-\allowbreak MiniLM-\allowbreak L12-\allowbreak v2}~\citep{reimers2020multilingual}, which
gives 384-d vectors. This covers both English (TIMME) and Slovak (Pokec), and it keeps
features comparable across datasets. Empty text maps to a zero vector, and
\texttt{has\_metadata} records whether text was present. On TIMME, the scores use
MiniLM and GloVe vectors concatenated. For content-aware spread, the scores use MiniLM
only and the receiver affinity uses GloVe only, so the spread model and the scorer
draw on different text sources (\cref{sec:method-block}).

\subsection{Parameters}
\Cref{tab:params} lists the pipeline parameters and \cref{tab:synthetic} the synthetic
grids.

\begin{table}[t]
\centering
\caption{Pipeline parameters. They were fixed before any blocking experiment was
run.}
\label{tab:params}
\small
\begin{tabularx}{\linewidth}{@{}l X@{}}
\toprule
Component & Setting \\
\midrule
Text encoder & \texttt{paraphrase-\allowbreak multilingual-\allowbreak MiniLM-\allowbreak L12-\allowbreak v2}, 384-d, text capped at 1000 characters \\
Topology view & Leiden, modularity objective, resolution 1.0, $R=100$ seeds (real data), 10 (synthetic) \\
Small communities & $<1\%$ of nodes $\to$ outside the label space ($t_i=-1$) \\
Content classifier & logistic regression, standardised inputs, $C=0.1$, 5 stratified folds \\
Label spreading & $\alpha = 0.5$, 40 iterations \\
Betweenness & sampled from 500 random source nodes \\
Spread models & weighted cascade; uniform $p = c\,p_c$, $c \in \{0.5, 1, 2, 4\}$ \\
Seeds per cascade & 10, drawn from nodes no strategy blocks \\
Node budgets (real data) & $k \in \{20, 50, 100, 200\}$ \\
Simulations & 100 (TIMME), 200 (Pokec slice) per spread model \\
Link budgets & TIMME $\{2000, 10000, 40000\}$; synthetic $\{100, 400, 1600\}$; 60 simulations, $c=2$ \\
\bottomrule
\end{tabularx}
\end{table}

\begin{table}[t]
\centering
\caption{Synthetic experiment grids.}
\label{tab:synthetic}
\small
\begin{tabularx}{\linewidth}{@{}l X X@{}}
\toprule
 & Detection benchmark & Regime map \\
\midrule
Nodes, communities & $n=4000$, $K=8$ & $n=4000$, $K=8$ \\
Degrees & power law (exponent 2.5), min.\ 8, cap 300 & same \\
Mixing $\mu$ & $\{0.1, 0.3, 0.5\}$ & $\{0.1, 0.3, 0.5\}$ \\
Text noise $\sigma$ & $\{0.5, 1, 2\}$ & 1 \\
Nuisance block & $\{0, 1.5\}$ & 1.5 \\
Mismatched share & 3\%, uniform & $\{1\%, 5\%, 10\%\}$, uniform or $\propto$ degree \\
Agreement effect $\beta$ & --- & $\{0, 1, 3, 6\}$ \\
Spread & --- & uniform IC, $c=2$; budget $n/100$; 60 simulations \\
Seeds per setting & 3 & 2 \\
Settings & 18 & 72 \\
Outcome & AUC, overall and within degree terciles & paired gain over best of PageRank / betweenness, 95\% CI \\
\bottomrule
\end{tabularx}
\end{table}


\subsection{Baselines}
\begin{itemize}
  \item \textbf{Detection}: consensus entropy, participation coefficient, content
        outlier and degree (\cref{sec:method-scores}). On TIMME there is also a
        shuffled-text version of each content score.
  \item \textbf{Node blocking}: no blocking, random, degree, adaptive degree,
        PageRank, sampled betweenness, Collective Influence, participation (alone and
        $\times$ PageRank), consensus entropy $\times$ PageRank, structural ambiguity
        $\times$ PageRank, content outlier $\times$ PageRank, and the structural gateway.
  \item \textbf{Link blocking}: random, degree product, PageRank product and the
        structural bridge score.
\end{itemize}

\subsection{Evaluation measures}
\begin{itemize}
  \item \textbf{Detection}: ROC AUC for recovering planted mismatches (synthetic) or
        party misfits (TIMME), overall and averaged within degree terciles. On TIMME
        we also report a one-sided Mann--Whitney $p$-value.
  \item \textbf{View quality}: modularity of each partition on the social graph,
        run-to-run stability (mean ARI between pairs of runs), and $\eta^2$ of log-degree
        across communities (does the partition just sort nodes by degree?). On TIMME
        we also report party NMI/ARI, but only as information.
  \item \textbf{Content informativeness}: accuracy of the cross-fitted text classifier
        in predicting the topology community, compared with always predicting the
        majority class.
  \item \textbf{Containment}: percentage reduction in infected share and in
        cross-community infected share compared with no blocking, at a fixed budget,
        with standard errors over simulations. Strategies are compared as paired
        differences under common random numbers.
\end{itemize}

\subsection{Implementation}
The pipeline is written in Python. It uses python-igraph (Leiden, PageRank,
betweenness), SciPy (sparse matrices, breadth-first reachability for live-edge
cascades), scikit-learn (classifiers, AUC, ARI/NMI) and sentence-transformers (text
encoding). The earlier GNN stage used PyTorch Geometric and ran on a rented CUDA GPU.
The current pipeline runs on CPU. Every module includes a self-check on a small planted
case with a known answer (for example, blocking the bridges of a two-community
barbell graph must cut cross-community spread more than blocking random nodes). Raw
outputs of every run are kept in the project's \texttt{results/} directory.
